\documentclass[../../main.tex]{subfiles}
\begin{document}
% SOURCE p=253; printed=238; object=ch08-translation-definition; region=upper
接下来研究平移在各种函数空间上的连续性。以下记号将在本章及下一章中使用。
\begin{definition}[平移与一致连续性]\label{definition:ch08-translation}
若 $f$ 是 $\mathbb R^n$ 上的函数、$y\in\mathbb R^n$，定义
\begin{align*}
\tau_yf(x)=f(x-y).
\end{align*}
若 $\|\tau_yf-f\|_u\to0$（$y\to0$），则称 $f$ \textbf{一致连续}。
\end{definition}
% SOURCE p=253; printed=238; object=ch08-translation-isometry; region=upper
注意，对 $1\leqslant p\leqslant\infty$ 有 $\|\tau_yf\|_p=\|f\|_p$，且 $\|\tau_yf\|_u=\|f\|_u$。上述一致连续性与通常的 $\varepsilon$--$\delta$ 定义等价，读者应核验这一点。
% SUPPLEMENT: id=v14-ch08-translation-details; parent=ch08-translation-isometry; kind=omitted-proof+motivation+hypotheses; origin=independent-derivation
\begin{remark}
\textbf{等距性、定义等价及动机（重构新增）。}对有限 $p$，Lebesgue 测度的平移不变性给

\begin{align*}
\int|f(x-y)|^p\,\mathrm dx=\int|f(u)|^p\,\mathrm du.
\end{align*}
$p=\infty$ 时，每个水平集只是作了平移，零测性不变，所以本质上确界相同。一致范数的等式来自 $x\mapsto x-y$ 为全空间的双射。平移保持零测集，故 $\tau_y$ 在 $L^p$ 的等价类上也定义良好。直接代入得 $\tau_y\tau_z=\tau_{y+z}$、$\tau_0=\mathrm{id}$，并有逆算子 $\tau_{-y}$。

若通常的一致连续性成立，给定 $\varepsilon>0$，用 $\varepsilon/2$ 选取 $\delta>0$，则 $|y|<\delta$ 时所有 $x$ 都有 $|f(x-y)-f(x)|<\varepsilon/2$，故一致范数至多 $\varepsilon/2<\varepsilon$。反向若平移差的一致范数趋零，对 $|u-v|<\delta$ 令 $y=v-u$、$x=v$，得到 $|f(u)-f(v)|\leqslant\|\tau_yf-f\|_u<\varepsilon$。所以二者等价。即使 $f$ 本身无界，差的上确界仍可用于这个定义；例如仿射函数一致连续。采用平移写法的动机，是将“邻近点的函数值接近”变成函数空间中一个算子族的连续性。负号是一致的记号约定：$\tau_y$ 把图像向 $y$ 的方向移。
\end{remark}
% END SUPPLEMENT: id=v14-ch08-translation-details
% SOURCE p=253; printed=238; object=ch08-compact-uniform-continuity; number=8.4; region=upper-middle
\begin{lemma}\label{lemma:ch08-compact-uniform-continuity}
若 $f\in C_c(\mathbb R^n)$，则 $f$ 一致连续。
\end{lemma}
\begin{proof}
给定 $\varepsilon>0$。对每个 $x\in\operatorname{supp} f$，由连续性可取 $\delta_x>0$，使 $|y|<\delta_x$ 时 $|f(x-y)-f(x)|<\varepsilon/2$。由于 $\operatorname{supp} f$ 紧，可取有限多个点 $x_1,\ldots,x_N$，使以 $x_j$ 为中心、半径为 $\delta_{x_j}/2$ 的球覆盖 $\operatorname{supp} f$。取 $\delta=\tfrac12\min_j\delta_{x_j}$，三角不等式给所有足够小的平移具有任意小的一致差，从而结论成立。具体的全空间估计如下。
\end{proof}
% SUPPLEMENT: id=v14-ch08-compact-uniform-details; parent=ch08-compact-uniform-continuity; kind=omitted-proof+motivation+hypotheses; origin=independent-derivation
\begin{remark}
\textbf{支撑内外的完整估计（重构新增）。}原证的逐点严格不等式取上确界后可能只留下非严格不等号，因此要得到 $\|\tau_yf-f\|_u<\varepsilon$，可在选 $\delta_x$ 时把上述预算 $\varepsilon/2$ 改为 $\varepsilon/4$。若支撑为空，则 $f=0$，无需选有限覆盖。否则使用这个预算取有限覆盖和 $\delta$。对任意 $u,v$ 满足 $|u-v|<\delta$，若两者都在支撑外，则 $f(u)=f(v)=0$。若至少一个在支撑内，称它为 $u$，选 $j$ 使 $|u-x_j|<\delta_{x_j}/2$，则
$|v-x_j|\leqslant|v-u|+|u-x_j|<\delta+\delta_{x_j}/2\leqslant\delta_{x_j}$。
因此
\begin{align*}
|f(u)-f(v)|\leqslant|f(u)-f(x_j)|+|f(x_j)-f(v)|<\varepsilon/2.
\end{align*}
取 $u=x$、$v=x-y$ 并取上确界，就得 $\|\tau_yf-f\|_u\leqslant\varepsilon/2<\varepsilon$。连续性提供各点的半径，紧性把它们变成一个正的统一半径，支撑外为零使估计覆盖整个非紧空间。紧支撑是充分条件，不是必要条件；有界一致连续函数可以不具紧支撑。
\end{remark}
% END SUPPLEMENT: id=v14-ch08-compact-uniform-details
% SOURCE p=253; printed=238; object=ch08-lp-translation-continuity; number=8.5; region=middle
\begin{proposition}\label{proposition:ch08-lp-translation-continuity}
若 $1\leqslant p<\infty$，则平移关于 $L^p$ 范数连续：对 $f\in L^p$、$z\in\mathbb R^n$，
\begin{align*}
\lim_{y\to0}\|\tau_{y+z}f-\tau_zf\|_p=0.
\end{align*}
\end{proposition}
\begin{proof}
由 $\tau_{y+z}=\tau_y\tau_z$，将 $f$ 换成 $\tau_zf$，只需考虑 $z=0$。先设 $g\in C_c$。当 $|y|\leqslant1$ 时，$\tau_yg$ 的支撑都在同一个紧集 $K$ 中。因此由前引理，
\begin{align*}
\int|\tau_yg-g|^p\leqslant\|\tau_yg-g\|_u^p m(K)\longrightarrow0
\qquad(y\to0).
\end{align*}
现在设 $f\in L^p$。给定 $\varepsilon>0$，由\cref{proposition:ch07-cc-lp-density}，存在 $g\in C_c$ 使 $\|g-f\|_p<\varepsilon/3$。于是
\begin{align*}
\|\tau_yf-f\|_p
&\leqslant\|\tau_y(f-g)\|_p+\|\tau_yg-g\|_p+\|g-f\|_p\\
&<2\varepsilon/3+\|\tau_yg-g\|_p.
\end{align*}
当 $y$ 足够小时最后一项小于 $\varepsilon/3$，故结论成立。
\end{proof}
% SUPPLEMENT: id=v14-ch08-lp-translation-details; parent=ch08-lp-translation-continuity; kind=omitted-proof+motivation+hypotheses; origin=independent-derivation
\begin{remark}
\textbf{共同支撑及条件作用（重构新增）。}可具体取 $K=\operatorname{supp} g+\overline B(0,1)$，它是两个紧集的乘积在加法映射下的连续像，因此紧且具有有限 Lebesgue 测度；$g=0$ 时直接处理。$\operatorname{supp}(\tau_yg)=\operatorname{supp} g+y\subseteq K$，而 $g$ 的支撑也在 $K$ 中，所以前一积分估计有根据。Lebesgue 测度为 Radon，故 $C_c$ 在有限 $p$ 的 $L^p$ 中稠密；平移等距性使逼近误差与 $y$ 无关。$p\geqslant1$ 用于范数的三角不等式，$p<\infty$ 用于上述稠密性及由有限体积上的一致误差推出 $L^p$ 误差。动机是先证明连续紧支撑函数的结论，再用统一有界的平移算子把它传给整个空间。对任意 $z$ 也可直接用等距性得到 $\|\tau_{y+z}f-\tau_zf\|_p=\|\tau_yf-f\|_p$。
\end{remark}
% END SUPPLEMENT: id=v14-ch08-lp-translation-details
% SOURCE p=253; printed=238; object=ch08-linfty-translation-failure; region=middle-lower
该命题对 $p=\infty$ 不成立；$L^\infty$ 范数与一致范数密切相关，因而这一点并不意外，见\cref{exercise:ch08-linfty-uc-representative}。
% SUPPLEMENT: id=v14-ch08-linfty-translation-counterexample; parent=ch08-linfty-translation-failure; kind=omitted-proof+motivation+hypotheses; origin=independent-derivation
\begin{remark}
\textbf{完整反例（重构新增）。}令 $f(x)=\chi_{\{x_1>0\}}$。对每个 $t>0$，$\tau_{te_1}f(x)=\chi_{\{x_1>t\}}$，故两函数在有正测度的条带 $0<x_1\leqslant t$ 上相差 $1$，且处处相差至多 $1$。因此 $\|\tau_{te_1}f-f\|_\infty=1$，即使 $t\downarrow0$ 也不趋零。有限 $p$ 下小集合上的误差可随集合体积缩小，而本质上确界只看误差高度；这解释了端点行为的差异。这个半空间示性函数属于 $L^\infty$，无需属于有限 $p$ 的空间。
\end{remark}
% END SUPPLEMENT: id=v14-ch08-linfty-translation-counterexample
% SOURCE p=253,254; printed=238-239; object=ch08-torus-definition; region=bottom-upper
\begin{definition}[周期函数与环面]\label{definition:ch08-torus}
我们的一些结论将涉及 $\mathbb R^n$ 上的多重周期函数。为简便起见，把各变量的基本周期统一取为 $1$：称 $f$ \textbf{周期}，是指对所有 $x\in\mathbb R^n$、$k\in\mathbb Z^n$ 均有 $f(x+k)=f(x)$。这样的函数完全由单位立方体
\begin{align*}
Q=[-\tfrac12,\tfrac12)^n
\end{align*}
上的值决定。周期函数可视为陪集空间 $\mathbb R^n/\mathbb Z^n\cong(\mathbb R/\mathbb Z)^n$ 上的函数；这个空间称为 \textbf{$n$ 维环面}，记作 $\mathbb T^n$。$n=1$ 时写作 $\mathbb T$。它是紧 Hausdorff 空间，通过映射
\begin{align*}
(x_1,\ldots,x_n)\longmapsto(e^{2\pi ix_1},\ldots,e^{2\pi ix_n})
\end{align*}
可与 $\{z\in\mathbb C^n:|z_j|=1\text{ 对所有 }j\}$ 识别。

另一方面，在测度论中把 $\mathbb T^n$ 与 $Q$ 识别；$\mathbb T^n$ 上的 Lebesgue 测度指由 $Q$ 上 Lebesgue 测度诱导的测度。特别地，$m(\mathbb T^n)=1$。环面上的函数既可看作 $\mathbb R^n$ 上的周期函数，也可看作 $Q$ 上的函数；必要时具体观点将由上下文明确。
\end{definition}
% SUPPLEMENT: id=v14-ch08-torus-details; parent=ch08-torus-definition; kind=omitted-proof+motivation+hypotheses; origin=independent-derivation
\begin{remark}
\textbf{商空间、基本域及测度的补证（重构新增）。}令 $q:\mathbb R^n\to\mathbb R^n/\mathbb Z^n$ 为商映射。每个 $x$ 唯一写成 $u+k$，其中 $u\in Q$、$k\in\mathbb Z^n$：逐坐标取 $k_j=\lfloor x_j+1/2\rfloor$ 即得存在性；同一半开长度为 $1$ 的区间中两个数若相差整数就必须相同，给出唯一性。半开条件因此避免边界重复，且不要求函数的最小正周期恰为 $1$。

记上述指数映射为 $E$。它连续满射到 $(S^1)^n$，且 $E(x)=E(y)$ 当且仅当 $x-y\in\mathbb Z^n$，所以诱导连续双射 $\widetilde E:\mathbb T^n\to(S^1)^n$。商空间是闭立方体 $[-1/2,1/2]^n$ 的连续像，故紧。紧空间到 Hausdorff 空间的连续双射为同胚：闭集紧，其像在 Hausdorff 空间中闭，故逆映射连续。因此 $\mathbb T^n$ 紧且 Hausdorff，并得到所述识别。逐坐标取陪集的映射也给 $\mathbb T^n\cong(\mathbb R/\mathbb Z)^n$，因为两者通过同一指数映射同胚于 $(S^1)^n$。

测度定义可写成 $m_{\mathbb T}(A)=m((q|_Q)^{-1}(A))$，定义域为其逆像 Lebesgue 可测的全部 $A$；$q|_Q$ 为双射，故这是一个完备概率测度。它与拓扑的 Borel 结构相容：$q|_Q$ 连续，而其逆也 Borel 可测。为证明后一项，令
\begin{align*}
Q_l=[-1/2,1/2-1/l]^n\quad(l\geqslant2),\qquad Q=\bigcup_{l=2}^\infty Q_l.
\end{align*}
则 $q|_{Q_l}$ 是紧集到 Hausdorff 空间的连续单射，因而是到其闭像上的同胚。任意 $Q$ 中 Borel 集 $B$ 的每个 $B\cap Q_l$ 的像均 Borel，故 $q(B)$ 为可数个 Borel 集之并。由 $m(Q)=1$ 得归一化。

动机是把周期边界粘合成紧空间，便于同时使用紧空间上的连续函数理论和有限测度空间上的积分理论。$Q$ 与环面在这里作集合及测度意义下的识别，不能据此声称半开立方体与环面拓扑同胚；后者紧而前者不紧。
\end{remark}
% END SUPPLEMENT: id=v14-ch08-torus-details
\subfile{section-01-exercises.tex}
\end{document}
